%
% template for Bachelor thesis
% LIACS
% January 7, 2019
%

\documentclass[12pt]{article}

% include some packages
\usepackage[left=2cm,right=2cm,top=2cm,bottom=3cm]{geometry}
\usepackage{graphicx}
\usepackage{tikz}
\usetikzlibrary{shapes,backgrounds,calc,arrows,arrows.meta}
\usepackage{microtype}
\usepackage{amsmath,amssymb}
\usepackage{booktabs}
\usepackage{float}
\usepackage{caption}
\usepackage{subcaption}
\usepackage{xcolor}

% --- Define the custom colors from the image ---
% Background colors (lighter)
\definecolor{bgorange}{RGB}{250, 235, 215}
\definecolor{bggreen}{RGB}{220, 255, 220}
\definecolor{bgblue}{RGB}{220, 220, 255}
\definecolor{bgred}{RGB}{255, 225, 225}

% Text and line colors (darker, more saturated)
\definecolor{textorange}{RGB}{238, 122, 44}
\definecolor{textgreen}{RGB}{89, 231, 89}
\definecolor{textblue}{RGB}{67, 83, 255}
\definecolor{textred}{RGB}{239, 65, 54}

% --- Helper command to wrap math equations in a TikZ node ---
% #1: Node name for referencing later
% #2: Background color
% #3: The math content
\newcommand{\mathbox}[3]{%
    \tikz[baseline=(#1.base), remember picture] {%
        \node[fill=#2, inner sep=3pt, anchor=base, rounded corners=2pt] (#1) {\ensuremath{\displaystyle #3}}%
    }%
}

% fill in title and author
\usepackage[setpagesize=false,colorlinks=true,linkcolor=blue,urlcolor=red,
  pdftitle={Evaluation of Causal Inference Methods for Gene Perturbation Prediction},
  pdfauthor={Luca Adrianus Verheul}]{hyperref}

\frenchspacing
\setlength{\parindent}{0pt}

\begin{document}

\thispagestyle{empty}


\includegraphics{logoleiden}

\vspace{-2.5cm}\hfill \begin{huge}\textbf{Opleiding Informatica}\end{huge}

\vspace{5cm}
\hfill \textbf{Evaluation of Causal Inference Methods for Gene Perturbation Prediction}

\vspace*{3mm}

\hfill Reconstructing and extending perturbation prediction metrics

\vspace*{14mm}

\begin{Large}
\hfill Luca Adrianus \textsc{Verheul}
\end{Large}

\vspace*{6.0cm}

\begin{large}

Supervisors:\\
Saber Salehkaleybar \& Gijs van Seeventer


\vspace*{2.8cm}
BACHELOR THESIS

\vspace*{5mm}
Leiden Institute of Advanced Computer Science (LIACS)\\
\href{www.liacs.leidenuniv.nl}{\underline{\texttt{www.liacs.leidenuniv.nl}}}\hfill 28/04/2025
\end{large}

\newpage



% abstract and references should fit on one page
\begin{abstract}
\noindent
Single-cell CRISPR perturbation screens enable systematic discovery of gene regulatory
interactions, but the perturbation space remains sparse: many genes (and conditions) are never
assayed experimentally. This motivates models that can \emph{predict unseen perturbation effects}
from limited interventional data.

In this thesis we evaluate \textbf{Bicycle} \cite{rohbeck2024bicycle}, a causal
model that combines a latent dynamical system with intervention-aware parameterization to infer a
\emph{cyclic} gene regulatory network and generalize to held-out perturbations. We reproduce the
Bicycle training and evaluation pipeline on three biological contexts from the Perturb-CITE-seq
study of Frangieh et al.\ \cite{frangieh2021perturbcite}, namely Control, IFN-$\gamma$, and
Co-culture.

Beyond the standard \emph{interventional mean absolute error} (I-MAE) used in the Bicycle paper,
we adopt two additional metrics from the \emph{Virtual Cell Challenge} \cite{virtualcellchallenge_eval}:
(1) the \emph{Differential Expression Score} (DES), measuring recovery of differential gene sets, and
(2) the \emph{Perturbation Discrimination Score} (PDS), measuring whether perturbation signatures
remain distinguishable even when effect sizes differ.

Our results show that I-MAE, DES, and PDS are correlated but not redundant: models with similar
I-MAE can differ substantially in DES/PDS, indicating that I-MAE alone may fail to reflect
biological correctness of predicted perturbation response signatures. We further observe a consistent model ranking across all metrics and discuss implications for evaluation practice in perturbation
modeling.
\end{abstract}

\bigskip

\thispagestyle{empty}
\tableofcontents
\thispagestyle{empty}

\clearpage
\setcounter{page}{1}

% -----------------------------------------------------------------------------
\section{Introduction}
Single-cell CRISPR perturbation technologies such as Perturb-seq have enabled systematic,
high-throughput interrogation of gene regulatory networks (GRNs). Instead of assaying one
perturbation at a time, these platforms simultaneously profile thousands of cells exposed to diverse
gene knock-outs or knock-downs, providing a uniquely detailed view into how gene expression
changes under intervention \cite{dixit2016perturbseq,adamson2016multiplexed,datlinger2017pooled}.
However, the perturbation space is always incomplete: many genes are
not perturbed in a given experiment, and testing all perturbations in all contexts is infeasible. This
motivates the need for computational models capable of \textbf{predicting the effect of unseen
perturbations} using only a subset of observed interventions.

The recently proposed \textbf{Bicycle} model (Intervention-Based Causal Discovery with Cycles)
addresses this challenge by combining causal inference with a latent dynamical systems model to
infer regulatory interactions directly from observational and interventional single-cell data
\cite{rohbeck2024bicycle}. Bicycle learns a \textbf{cyclic} causal graph and uses it to generalize across
perturbations, enabling predictions for conditions not encountered during training. Cycles are
particularly important in biology, where feedback regulation is common (e.g., signaling loops and
autoregulation).

\subsection{Thesis scope and research questions}
This thesis evaluates Bicycle on a subset of the Perturb-CITE-seq dataset of
Frangieh et al.\ \cite{frangieh2021perturbcite}, extending the original evaluation by incorporating
three complementary metrics used in the \textbf{Virtual Cell Challenge}:
\begin{itemize}
    \item the \emph{Interventional Mean Absolute Error} (I-MAE) already considered in Bicycle,
    \item the \emph{Differential Expression Score} (DES),
    \item the \emph{Perturbation Discrimination Score} (PDS).
\end{itemize}
These metrics assess distinct aspects of model performance, and let us test whether I-MAE alone is
a reliable indicator of biologically meaningful prediction quality, definitions are provided in Section~3.

Concretely, we address the following questions:
\begin{enumerate}
    \item \textbf{Reproducibility:} Can we reproduce the quantitative behavior and ranking trends
    reported for Bicycle on the Frangieh data, and what deviations arise in absolute I-MAE?
    \item \textbf{Metric agreement:} How strongly do I-MAE, DES, and PDS correlate across
    contexts and training runs, and where do they disagree?
    \item \textbf{Metric validity:} In what ways can I-MAE be misleading as a \emph{standalone}
    performance metric for perturbation prediction?
\end{enumerate}

\subsection{Datasets}
We evaluate Bicycle across three biologically distinct contexts from the Frangieh dataset:
\begin{itemize}
    \item \textbf{Control}: baseline T-cell culture without cytokine stimulation or without tumor-cell co-culture.
    \item \textbf{IFN-}\(\gamma\): interferon-$\gamma$ stimulation, inducing an immune-activated transcriptional state and altering baseline pathway activity.
    \item \textbf{Co-culture}: T cells cultured with melanoma cells, introducing tumor--immune interaction signals and additional environmental variation.
\end{itemize}
These contexts differ in baseline pathway activation and regulatory state, so the same genetic perturbation can yield context-dependent downstream expression changes (e.g., IFN-$\gamma$ activates immune pathways and can shift regulatory sensitivity, while co-culture introduces tumor--immune signaling cues).

\subsection{Contributions}
The main contributions of this thesis are:
\begin{itemize}
    \item A reconstruction of the Bicycle pipeline for the Frangieh contexts with careful reporting
    of preprocessing, training, and evaluation decisions.
    \item An extended evaluation that combines I-MAE with two Virtual Cell Challenge metrics
    (DES, PDS) to provide a multidimensional view of performance.
    \item An empirical analysis and discussion of why I-MAE may not be a sufficient proxy for
    biologically meaningful perturbation prediction quality.
\end{itemize}

% -----------------------------------------------------------------------------
\section{Background \& Preliminaries}
This section discusses the methodology of the Bicycle model based directly on the original
paper \cite{rohbeck2024bicycle}. We also introduce the causal and statistical concepts needed
for interpreting Bicycle’s training objective and evaluation.

\subsection{Single-cell perturbation screens and notation}
We consider a dataset of cells indexed by \(n \in \{1,\dots,N\}\). Each cell is measured across \(G\)
selected genes. Let \(x_n \in \mathbb{N}^G\) denote the observed unique molecular identifier (UMI) counts for cell \(n\). Cells are
grouped into \emph{contexts} (Control / IFN-\(\gamma\) / Co-culture) and into \emph{perturbation
conditions}. A perturbation condition \(i\) targets a set of genes \(P^i \subseteq \{1,\dots,G\}\).

Throughout, we distinguish (i) \emph{directly perturbed} genes \(g \in P^i\) and (ii) \emph{non-target}
(or \emph{downstream}) genes \(g \notin P^i\). Bicycle’s interventional evaluation focuses on the
non-target set because predicting downstream consequences is the goal.

\subsection{Causal graphs and interventions}
Bicycle is motivated by \emph{structural causal models} (SCMs). An SCM specifies variables
\(X_1,\dots,X_G\) through \emph{structural assignments}
\[
X_g := f_g\!\big(X_{\mathrm{Pa}(g)}, \varepsilon_g\big),
\]
where \(\mathrm{Pa}(g)\) denotes the \emph{parents} (direct causes) of \(X_g\) and \(\varepsilon_g\) is exogenous
noise. The associated directed graph has an edge \(h \to g\) whenever \(X_h\) appears in \(f_g\). In our
setting, variables correspond to gene expression levels, so expression of gene \(g\) is generated as a
function of its regulators and noise. \emph{Interventions} correspond to replacing the structural
assignment(s) of targeted genes while leaving the remaining mechanisms unchanged. Classical causal
discovery algorithms often assume acyclicity, i.e., the causal graph is a DAG. However, biological GRNs
contain feedback loops, motivating cyclic causal models
\cite{mooij2013cyclic,bongers2021theoretical}.

\begin{figure}[H]
\centering
\includegraphics[width=0.7\linewidth]{paper/figures/bicycle-figure-1.png}
\caption{Bicycle integrates information from multiple interventions to infer an underlying causal graph and predict responses for unseen interventions \cite{rohbeck2024bicycle}.}
\label{fig:bicycle_unseen_interventions}
\end{figure}

\subsection{Latent Dynamical System}
Let \(z \in \mathbb{R}^G\) denote the latent, noise-free gene expression state. Bicycle assumes that gene
expression follows a multivariate Ornstein--Uhlenbeck (OU) process:
\[
dz_g(t) =
\left(
\alpha_g + \sum_{g \neq h} \beta_{hg} z_h(t) - z_g(t)
\right) dt
+ \sum_h \sigma_{gh}\, dW_h(t),
\]
where \(\alpha_g\) is the basal transcription rate, \(\beta_{hg}\) is the direct causal effect of gene \(h\) on
\(g\), and \(W_h(t)\) are Wiener processes. Define the drift matrix \(B = I_G - \beta^\top\), where
\(I_G \in \mathbb{R}^{G \times G}\) is the identity matrix and \(\beta^\top\) is the transpose of the regulatory
coefficient matrix \(\beta \in \mathbb{R}^{G \times G}\) with entries \((\beta)_{hg} = \beta_{hg}\).
When \(B\) has eigenvalues with positive real parts, this system admits a closed-form steady-state distribution \(p_{\mathrm{ss}}(z)\), i.e., the distribution approached as \(t \to \infty\)
\cite{wang1945steadystate}:
\[
p_{\mathrm{ss}}(z) = \mathcal{N}(\bar{z}, \omega)
\quad\text{with}\quad
\bar{z} = B^{-1}\alpha,
\]
and covariance \(\omega\) given by the Lyapunov equation:
\[
B\,\omega + \omega B^\top = \sigma\sigma^\top.
\]
This formulation links continuous-time dynamics to cyclic causal structure.

\subsection{Observation Model}
Observed single-cell RNA-seq counts are \emph{compositional}: for a given cell \(n\), sequencing produces a finite
total number of captured molecules (UMIs), so only relative abundances across genes are directly
observable. Bicycle therefore models the observed counts \(x_n \in \mathbb{N}^G\) using a Multinomial
likelihood,
\[
x_n \sim \mathrm{Multinomial}(N_n,\pi_n),
\qquad
\pi_n = \mathrm{Softmax}(z_n),
\]
where \(N_n = \sum_{g=1}^G x_{ng}\) is the (observed) library size and \(z_n \in \mathbb{R}^G\) denotes the
latent log-proportion state. In this view, \(z_n\) captures the underlying biological state, while the
Multinomial distribution captures sampling variability induced by finite sequencing depth.

\subsection{Perturbation Modeling via Independent Causal Mechanisms}
A key innovation of Bicycle is incorporating perturbation-specific modifications to the dynamical
system parameters. For a perturbation condition \(i\) targeting genes \(P^i\), Bicycle defines
perturbation-specific parameters
\[
\alpha^i_g, \quad
\beta^i_{hg}, \quad
\sigma^i_{gh},
\]
but \textbf{only} for genes directly targeted by the intervention. For all other genes
\(g \notin P^i\), parameters remain unchanged:
\[
\alpha^i_g = \alpha_g,\quad
\beta^i_{hg} = \beta_{hg},\quad
\sigma^i_{gh} = \sigma_{gh}.
\]
This realizes the \textbf{Independent Causal Mechanisms} (ICM) principle, which posits that
different causal modules remain stable across environments unless directly intervened upon
\cite{scholkopf2021causalrep}. In Bicycle, ICM becomes an explicit inductive bias: changes are
localized to the perturbed gene’s mechanism.

\subsection{Variational Inference}
Bicycle optimizes a variational Evidence Lower Bound (ELBO) with additional structural penalties \cite{rohbeck2024bicycle}.
Let \(q(z_n)\) be a factorized Gaussian variational distribution. The objective is:

\vspace*{2cm}

\begin{equation}
\begin{split}
\mathcal{L}^{\gamma,\xi,\lambda}(x, \alpha, \beta, \sigma, \hat{\alpha}, \hat{\sigma}, \mu_z, \sigma_z) &= \\
\frac{1}{N} \sum_{n=1}^{N} \bigg[ \mathbox{term1}{bgorange}{\langle \log p(x_n \mid z_n) \rangle_{q(z_n|\mu_{z_n},\sigma_{z_n})}} &- \gamma \mathrm{KL}\big(q(z \mid \mu_{z_n}, \sigma_{z_n}) \,\|\, \mathbox{term2}{bggreen}{\mathcal{N}(z; \bar{z}^{i_n}, \omega^{i_n}) \big)} \bigg] \\
&- \frac{\xi}{I} \sum_{i=0}^{I} \mathbox{term3}{bgblue}{\big\| B^i \omega^i + \omega^i (B^i)^T - \sigma^i (\sigma^i)^T \big\|_2^2} - \lambda \mathbox{term4}{bgred}{\|\beta\|_1}
\end{split}
\end{equation}

\begin{tikzpicture}[overlay, remember picture, >=Latex, font=\sffamily\footnotesize, thick]

    % --- Orange Annotation ---
    \draw[textorange, <-] (term1.south) -- ++(0, -0.8cm) -| ($(term1.west) + (2cm, -1.2cm)$) coordinate (o1);
    \node[textorange, anchor=north] at (o1) {Log-Likelihood of observed counts};

    % --- Green Annotation ---
    \draw[textgreen, <-] (term2.north) -- ++(0, 1.5cm) coordinate (g1);
    \node[textgreen, anchor=south west] at (g1 -| term1.west) {Steady-State distribution of latent expression};
    \draw[textgreen] (g1) -- (g1 -| term1.west);

    % --- Blue Annotation ---
    \draw[textblue, <-] (term3.south) -- ++(0, -1.5cm) coordinate (b1);
    \node[textblue, anchor=south east] at (b1 -| term3.west) {$\Delta$ between LHS and RHS of Lyapunov Eq.};
    \draw[textblue] (b1) -- (b1 -| term3.west);

    % --- Red Annotation ---
    \draw[textred, <-] (term4.south) |- ++(0, -0.8cm) coordinate (r1);
    \node[textred, anchor=east] at (r1 -| term4.west) {Sparsity};
    \draw[textred] (r1) -- (r1 -| term4.west);

\end{tikzpicture}

\subsection{How I-MAE is defined and why it is graph-dependent}
\label{sec:imae-graph}
Bicycle evaluates generalization using \emph{Interventional Mean Absolute Error} (I-MAE, defined in Section~\ref{sec:imae-def}) on held-out perturbations. The key point is
that Bicycle does \emph{not} simply output a mean expression vector per perturbation. Instead, the
metric evaluates how well the model predicts each gene \(d\) in each cell, \emph{conditioned on all
other genes} \(\neg d\) in that cell \cite{rohbeck2024bicycle}.

In words, Bicycle approximates the conditional distribution \(p(x_{nd} \mid x_{n\neg d})\) via a
variational approximation over the latent variables. The I-MAE is then computed on \emph{non-target}
genes \(d \in G\setminus P^i\) for each held-out intervention \(i\). This makes I-MAE \emph{graph-dependent} because the learned regulatory structure determines which conditional dependencies the model uses when predicting each gene from the others; it is therefore not merely an unconditional ``mean profile'' error.

\subsection{Reconstructing Bicycle and reproducibility considerations}
In addition to implementing Bicycle as described in \cite{rohbeck2024bicycle}, we attempted to
\emph{closely reproduce} the I-MAE values reported in the original paper.

Across repeated runs, we matched the hyperparameters, training schedule, loss-weighting
coefficients, sparsity penalty grid for \(\lambda\), and preprocessing choices to the extent they were
documented in \cite{rohbeck2024bicycle} and the official code. We also matched the reported gene
subset (\(G=61\)) and the 90/10 perturbation split protocol (55 train perturbations, 6 test).

Overall, we were able to closely recreate the reported I-MAE behavior (including qualitative
rankings across contexts and the spread across runs), but we did not obtain a perfect quantitative
match to the absolute I-MAE values reported by \cite{rohbeck2024bicycle}. This is plausible given
that several training details of the original Perturb-seq runs (e.g., epoch budget and optimizer
settings) are not fully specified in the paper, and because our reproduction differs in a few
implementation choices. Concretely, we trained for up to 40{,}000 epochs (plus 5{,}000 epochs of
latent pretraining) with early stopping, used Adam with a learning rate of $10^{-3}$ and batch size
1024, explored a wider regularization sweep (including spectral and Lyapunov scaling), and used a
full-rank covariance factorization (rank $=61$) for numerical stability, whereas the paper motivates
a low-rank covariance parameterization for large real-world datasets.
In our experiments, I-MAE values tended to fall near the lower end of the
ranges shown in the paper’s boxplots, and this discrepancy persisted across contexts and random
seeds.

This gap is plausible because I-MAE is sensitive to small details of (i) conditional inference,
(ii) how control cells are handled, (iii) normalization/log-transform choices applied after
predicting compositional probabilities, and (iv) the number of runs. Minor differences that are conceptually equivalent can still
shift absolute I-MAE. Importantly, the \emph{relative} behavior of Bicycle---context difficulty ranking
and generalization trends---was consistent with \cite{rohbeck2024bicycle}, so the reconstructed
pipeline remains a valid foundation for the extended evaluation with I-MAE, DES, and PDS.

% -----------------------------------------------------------------------------
\section{Motivation}
Most prior work evaluates perturbation prediction using a reconstruction error metric such as the
Interventional Mean Absolute Error (I-MAE). While simple and interpretable, I-MAE alone does not
capture biologically meaningful structure.

In the \textbf{Virtual Cell Challenge}, this motivated the introduction of two additional metrics---the
Differential Expression Score (DES) and the Perturbation Discrimination Score (PDS)---designed to
evaluate perturbation modeling beyond reconstruction error \cite{virtualcellchallenge_eval}.

\subsection{Interventional Mean Absolute Error (I-MAE)}
\label{sec:imae-def}
The I-MAE is defined as the average absolute error between predicted and true expression under
held-out perturbations, evaluated on \emph{non-target} genes. In Bicycle, predictions are obtained
from a conditional distribution and converted into normalized log-expression values
\cite{rohbeck2024bicycle}.

Although useful for quantifying reconstruction accuracy, I-MAE has several limitations as a \emph{sole}
metric:
\begin{itemize}
    \item \textbf{Many genes do not change:} If only a small subset of genes responds to the perturbation,
    a model can obtain a low average error by predicting the many unaffected genes well, even if it misses
    the key responders.

    \item \textbf{Averaging can hide wrong biology:} Two predictions can have similar MAE even if one gets the
    direction/pattern of the perturbation response wrong but compensates by being close to the mean.

    \item \textbf{No notion of distinguishability:} MAE does not test whether different perturbations produce
    distinct predicted response signatures; distinct perturbations can collapse toward similar predicted deltas
    while still keeping average error moderate.
\end{itemize}

\subsection{Differential Expression Score (DES)}
Differential expression (DE) refers to genes whose expression differs significantly between perturbed and control cells under a statistical test.
DES evaluates whether the model correctly predicts \textbf{which genes become significantly
differentially expressed} after a perturbation. For each perturbation \(k\), predicted and true DE
genes are identified by comparing perturbed cells to controls (Wilcoxon test with multiple-testing
correction and an FDR threshold as in the Virtual Cell Challenge definition) \cite{virtualcellchallenge_eval}. Here the Wilcoxon rank-sum test is a nonparametric test comparing expression distributions between perturbed and control groups, and FDR is the false discovery rate controlled via multiple-testing correction across genes. Here \(k\) indexes perturbations (i.e., target genes) in the evaluation set; \(G_{k,\mathrm{true}}\) and
\(G_{k,\mathrm{pred}}\) denote the differentially expressed (DE) gene sets for perturbation \(k\). The
per-perturbation score is
\[
\mathrm{DES}_k = \frac{\lvert G_{k,\mathrm{pred}} \cap G_{k,\mathrm{true}} \rvert}{\lvert G_{k,\mathrm{true}} \rvert}.
\]
To prevent trivial solutions where a model predicts ``everything is DE'', the evaluation constrains
predicted DE sets to match the size of the true DE set \cite{virtualcellchallenge_eval}. DES
therefore focuses on \emph{which} genes change rather than absolute magnitude across the full
profile.

\subsection{Perturbation Discrimination Score (PDS)}
PDS assesses whether each perturbation’s signature is \textbf{recognizably distinct}. In the Virtual Cell
Challenge rubric \cite{virtualcellchallenge_eval}, the score is computed by:
\begin{enumerate}
    \item \textbf{Compute pseudobulk profiles.} For each perturbation \(k \in \{1,\dots,N\}\), compute the
    true pseudobulk expression profile \(x_k\) by averaging the \(\log(1+x)\)-normalized expression of each
    gene across all cells under perturbation \(k\), and compute the corresponding non-targeting control
    pseudobulk \(x_{\mathrm{ntc}}\). Analogously, compute predicted pseudobulk profiles \(\hat{x}_k\) and
    \(\hat{x}_{\mathrm{ntc}}\).

    \item \textbf{Form delta signatures.} Define true and predicted delta signatures as
    \[
        \delta_k = x_k - x_{\mathrm{ntc}},
        \qquad
        \hat{\delta}_k = \hat{x}_k - \hat{x}_{\mathrm{ntc}}.
    \]

    \item \textbf{Rank by L1 similarity and score.} For each predicted perturbation \(p\), compute the
    L1 (Manhattan) distance to every true delta signature,
    \[
        d_{pt} = \|\hat{\delta}_p - \delta_t\|_1, \qquad t=1,\dots,N,
    \]
    sort \(\{d_{pt}\}_{t=1}^N\) in ascending order, and let \(r_p\) be the rank position of the correct
    match \(t=p\). The per-perturbation score is then
    \[
        \mathrm{PDS}_p = 1 - \frac{r_p - 1}{N}.
    \]
\end{enumerate}
The rank is normalized to \([0,1]\) such that higher is better. PDS is designed to be robust to global
effect-size scaling: if the correct \emph{pattern} is present, a model may score well even if it
underestimates magnitude.

\subsection{Why All Three Metrics Are Needed}
The three performance metrics capture complementary information:
\begin{itemize}
    \item \textbf{I-MAE}: overall conditional reconstruction accuracy,
    \item \textbf{DES}: whether the correct DE genes were identified,
    \item \textbf{PDS}: whether perturbations remain distinguishable.
\end{itemize}
A model may have low I-MAE but fail to capture the correct biological response (low DES), or may
capture the response structure but not its magnitude (high PDS, moderate I-MAE). Evaluating all
three metrics provides a multidimensional assessment of perturbation prediction quality.

% -----------------------------------------------------------------------------
\section{Related Work}\label{relatedwork}
We position Bicycle and our evaluation approach in the broader literature on (i) single-cell
perturbation screens, (ii) perturbation response prediction, (iii) causal discovery with cycles, and
(iv) benchmarking and metric design.

\subsection{Single-cell perturbation datasets}
Early pooled CRISPR + scRNA-seq methods such as Perturb-seq enabled scalable mapping from
genetic interventions to transcriptional states \cite{dixit2016perturbseq,adamson2016multiplexed,jaitin2016immune}.
Subsequent improvements increased throughput and expanded modalities (e.g., CITE-seq proteins)
\cite{mimitou2019multimodal,datlinger2021ultra}.

Frangieh et al.\ introduced \emph{Perturb-CITE-seq} to profile both RNA and surface proteins in
patient-derived melanoma--TIL co-cultures \cite{frangieh2021perturbcite}. A key feature of this study is
explicit \emph{context dependence}: the same genetic perturbation can induce different transcriptional programs
depending on immune pressure and culture condition (Control vs IFN-$\gamma$ vs Co-culture). This makes the
dataset a challenging and realistic testbed for out-of-distribution (OOD) perturbation prediction.

\subsection{Perturbation response prediction models}
A major line of work focuses on predicting expression changes under unseen perturbations without
explicitly recovering a causal GRN. Representative approaches include variational autoencoder (VAE) models such as scGen \cite{lotfollahi2019scgen}
(and extensions for compositional/complex settings such as CPA \cite{lotfollahi2023cpa}),
conditional generative models such as PerturbNet \cite{yu2025perturbnet},
graph-based deep learning approaches for single- and multi-gene perturbations such as GEARS \cite{roohani2024gears},
and causal meta-learning approaches such as CaML \cite{nilforoshan2023zeroshot}.

Bicycle differs by aiming to infer an \emph{interpretable} regulatory matrix \(\beta\) grounded in a
steady-state dynamical system and by encoding interventions through ICM-style parameter
invariance.

\subsection{Causal discovery with cycles and interventions}
Many causal discovery methods assume acyclicity, including among others NODAGS and NOTEARS \cite{sethuraman2023nodagsflow,zheng2018notears}.
Interventional extensions exist, but explicit acyclicity constraints preclude modeling feedback.
Cyclic causal discovery has been studied in equilibrium and dynamical settings
\cite{mooij2013cyclic,bongers2021theoretical,sethuraman2025dccdconf}. Bicycle contributes by combining cyclic structure
with intervention-specific mechanism changes and a count-based likelihood appropriate for scRNA-seq.

\subsection{Benchmarks and evaluation metrics}
Benchmarking perturbation predictors is non-trivial: different metrics reward different aspects of
performance. The Virtual Cell Challenge proposes an evaluation framework aimed at context
generalization and uses DES and PDS alongside MAE \cite{virtualcellchallenge_eval}. This thesis follows
that philosophy: I-MAE is informative but insufficient as a single ``leaderboard metric''.

Related work suggests two recurring themes:
\begin{itemize}
    \item \textbf{Causal structure matters} for out-of-distribution generalization under intervention, particularly in biological systems with feedback.
    \item \textbf{Metric choice matters} because different metrics can invert model rankings.
\end{itemize}
Our experiments therefore emphasize multi-metric evaluation.

% -----------------------------------------------------------------------------
\section{Experiments}\label{experiments}

\subsection{Dataset and splits}
We use the Frangieh Perturb-CITE-seq dataset as processed in the Bicycle paper
\cite{rohbeck2024bicycle,frangieh2021perturbcite}. Following \cite{rohbeck2024bicycle}, we
restrict to the same subset of \(G=\texttt{61}\) genes and evaluate each biological context separately. In this processed subset, each perturbation corresponds to a single target gene from the \(G=61\)-gene panel.
For readability, we refer to perturbations by their target-gene index (1--61). The train/test split is therefore
a split over perturbations (targets), not over measured genes in the expression vector.

\textbf{Perturbation split.} We follow the reported protocol: 90\% of perturbations are used for
training, 10\% are held out for testing. The held-out set is unseen during training.

\textbf{Controls.} Control (non-targeting) cells are used to define baseline expression and are
necessary for computing delta signatures (PDS) and DE testing (DES). In our pipeline we ensure
that each context has an explicit control subset aligned with the perturbation condition metadata. Table \ref{tab:dataset-stats} summarizes the resulting dataset sizes and split counts per context.

\begin{table}[H]
\centering
\caption{Dataset statistics used in our experiments.}
\label{tab:dataset-stats}
\begin{tabular}{lrrrrrr}
\toprule
Context & \#Cells & \#Control & \#Perturbed & \#Perts & \#Train & \#Test \\
\midrule
Control & 21{,}155 & 13{,}142 & 8{,}013  & 61 & 55 & 6 \\
IFN-$\gamma$ & 36{,}312 & 23{,}771 & 12{,}541 & 61 & 55 & 6 \\
Co-culture & 31{,}694 & 19{,}667 & 12{,}027 & 61 & 55 & 6 \\
\bottomrule
\end{tabular}
\end{table}

\subsection{Preprocessing}
We match the Bicycle paper’s preprocessing as closely as possible.
\begin{itemize}
    \item \textbf{Bicycle training:} trained on raw UMI counts with a Multinomial likelihood.
    \item \textbf{Metric computation:} For DES/PDS computation we follow the Virtual Cell Challenge metric definitions, which operate on normalized log-expression and pseudobulk delta signatures based on \(\log(1+x)\)-transformed, library-size-normalized counts.
\end{itemize}

\subsection{Model training and hyperparameters}
We train Bicycle with the hyperparameters described in \cite{rohbeck2024bicycle}.
\begin{itemize}
    \item \(\gamma = 1.0\), \(\xi = 1.0\).
    \item \(\lambda \in \{0.01, 0.1, 1.0\}\) tuned on training perturbations.
    \item Optimizer: Adam.
    \item Batch size \(=1024\), learning rate \(=10^{-3}\), number of epochs \(=40000\) (plus latent pretraining for \(5000\) epochs), early stopping enabled with patience \(=500\) and \(\Delta_{\min}=0.01\); number of seeds/runs: 5.
\end{itemize}

\subsection{Evaluation protocol}
We evaluate on held-out perturbations for each context and report performance across repeated training seeds.

\subsubsection{Metric 1: I-MAE}
We compute I-MAE as in the Bicycle paper: for each held-out perturbation \(i\) and each cell \(n\),
we evaluate a conditional distribution for each non-target gene \(d\in G\setminus P^i\), convert to
expected normalized log-expression, and compute MAE.

\subsubsection{Metric 2: DES}
For each perturbation \(k\), we compare perturbed cells to non-targeting controls using a Wilcoxon
rank-sum test per gene, apply multiple-testing correction, define the true/predicted DE gene
sets at the chosen FDR, and compute overlap as in the Virtual Cell Challenge rubric
\cite{virtualcellchallenge_eval}.

\subsubsection{Metric 3: PDS}
We compute pseudobulk expression per perturbation and controls, compute delta signatures, and
rank predicted deltas by similarity to true deltas using Manhattan (\(\ell_1\)) distance; the normalized
rank yields PDS \cite{virtualcellchallenge_eval}.

\subsubsection{Baselines for I-MAE}
\label{sec:baselines}
In addition to \textsc{Bicycle}, we report two point-prediction baselines. We evaluate all methods under the same I-MAE functional:
for each perturbed cell \(n\) under intervention \(i\), and for each non-target gene
\(d \in \mathcal{G}\setminus \mathcal{P}_i\), we compare the true normalized log-expression \(y_{nd}\)
to a predicted value \(\hat{y}_{nd}\).
For \textsc{Bicycle}, \(\hat{y}_{nd}\) is obtained from the model-induced conditional expectation
(as described in Section~\ref{sec:imae-graph}).
For baselines, \(\hat{y}_{nd}\) is defined directly as a fixed reference profile (i.e., we bypass conditional inference
but still compute the same per-gene absolute error).

\begin{itemize}
  \item \textbf{WT-Baseline (``no perturbation effect'').}
  \[
    \hat{y}_{nd} := \mu^{\mathrm{ctrl}}_{d},
  \]
  where \(\mu^{\mathrm{ctrl}}_{d}\) is the mean normalized log-expression of gene \(d\) over non-targeting control cells
  in the same biological context.
  This tests whether any method improves over predicting the control state for every perturbed cell.

  \item \textbf{VCC cell-mean baseline (``global mean'').}
  \[
    \hat{y}_{nd} := \mu^{\mathrm{train}}_{d},
  \]
  where \(\mu^{\mathrm{train}}_{d}\) is the mean normalized log-expression of gene \(d\) over all training cells
  (controls + training perturbations) in the same biological context.
  This tests whether models outperform a context-specific global mean, analogous in spirit to the Virtual Cell Challenge
  cell-mean reference.
\end{itemize}

Because these baselines ignore perturbation identity, they can score deceptively well on reconstruction-style metrics when
perturbation effects are weak or sparse, motivating complementary metrics such as DES and PDS.

\subsubsection{Metric agreement analysis}
To quantify agreement between metrics across contexts and repeated runs, we compute pairwise correlations using both Pearson's \(r\) (linear association) and Spearman's \(\rho\) (rank association). Since I-MAE is a loss (lower is better) while DES and PDS are scores (higher is better), negative correlations between I-MAE and DES/PDS are expected.

\subsection{Robustness and sensitivity analysis}\label{robustness}
DES and PDS depend on design choices (DE test, FDR threshold, similarity metric, and scaling).
While we follow the Virtual Cell Challenge rubric, we make these choices explicit:
(i) DE testing uses a two-sided Wilcoxon rank-sum test with tie correction;
(ii) multiple-testing correction is Benjamini--Hochberg with an FDR threshold of \(0.05\);
(iii) PDS ranks delta signatures by Manhattan (\(\ell_1\)) distance; and
(iv) scaling uses per-cell library-size normalization followed by a \(\log(1+x)\) transform.
Concretely, for raw counts \(x_{ng}\) we compute \(L_n=\sum_{g=1}^{G}x_{ng}\), normalize
\(\tilde{x}_{ng}=x_{ng}/L_n\), and use \(y_{ng}=\log(1+\tilde{x}_{ng})\) for DE testing and
pseudobulk/delta computations.

\subsubsection{DES sensitivity}
DES depends on how the ``true'' DE set \(G_{k,\mathrm{true}}\) is defined by the DE-testing pipeline
(test statistic, multiple-testing correction, and FDR threshold). Changing the FDR threshold changes
\(\lvert G_{k,\mathrm{true}} \rvert\) and therefore changes the DES values. We therefore interpret DES as a
metric that is meaningful \emph{relative to a fixed DE-calling procedure} (in our case, Wilcoxon + BH at
FDR \(=0.05\)), rather than as an absolute quantity comparable across different DE pipelines.

\subsubsection{PDS sensitivity}
PDS depends on the similarity measure used to rank delta signatures and on whether signatures are
scaled or standardized before similarity computation. Recent work notes that distance metrics and
scaling can shift PDS rankings across methods \cite{liu2025pds}. For this thesis, we therefore treat
PDS as a complement to (not a replacement for) I-MAE and DES, emphasizing agreement patterns
and rank stability rather than absolute values.

\subsection{Main results}

\subsubsection{Aggregate I-MAE performance by context}

\begin{figure}[H]
    \centering
    \includegraphics[width=\linewidth]{paper/figures/imae_boxplots.pdf}
    \caption{Comparison of out-of-distribution generalisation in terms of the interventional mean absolute error (I-MAE) across three biological contexts. Co-culture: lymphocytes co-cultured with melanoma cells. IFN-$\gamma$: lymphocytes stimulated by interferon-$\gamma$. Control: lymphocytes cultured in neutral medium. Methods shown are \textit{VCC-Baseline} (the Virtual Cell Challenge cell-mean baseline as defined in Section~\ref{sec:baselines}), \textit{WT-Baseline} (the control-mean baseline as defined in Section~\ref{sec:baselines}), and \textit{Bicycle}. Box plots show the distribution of I-MAE over repeated runs (minimum, lower quartile, median, upper quartile, maximum, and outliers).}
    \label{fig:imae_boxplots}
\end{figure}

\begin{table}[H]
\centering
\caption{Aggregate performance across runs (mean $\pm$ std). Test aggregates over the 6 held-out perturbations
(target genes 56--61, 1-indexed).}
\label{tab:main-metrics-1}
\begin{tabular}{lccccc}
\toprule
Context
& I-MAE \(\downarrow\)
& I-MAE (WT) \(\downarrow\)
& I-MAE (VCC) \(\downarrow\)
& DES \(\uparrow\)
& PDS \(\uparrow\)\\
\midrule
Control
& 0.723 $\pm$ 0.080
& 0.797
& 0.79
& 0.969 $\pm$ 0.000
& 0.389 $\pm$ 0.007\\
IFN-$\gamma$
& 0.874 $\pm$ 0.027
& 0.955
& 0.94
& 0.938 $\pm$ 0.000
& 0.374 $\pm$ 0.004\\
Co-culture
& 0.719 $\pm$ 0.014
& 0.758
& 0.753
& 0.954 $\pm$ 0.000
& 0.590 $\pm$ 0.000\\
\bottomrule
\end{tabular}
\end{table}

\begin{table}[H]
\centering
\caption{Aggregate performance across runs (mean $\pm$ std). Train aggregates over the 55 training perturbations
(target genes 1--55, 1-indexed).}
\label{tab:main-metrics-2}
\begin{tabular}{lcc}
\toprule
Context & DES \(\uparrow\) & PDS \(\uparrow\)\\
\midrule
Control      & 0.962 $\pm$ 0.002 & 0.558 $\pm$ 0.001\\
IFN-$\gamma$ & 0.970 $\pm$ 0.003 & 0.540 $\pm$ 0.002\\
Co-culture   & 0.985 $\pm$ 0.000 & 0.546 $\pm$ 0.000\\
\bottomrule
\end{tabular}
\end{table}

\textbf{Dataset difficulty ranking.} On the held-out test interventions (the 6 unseen perturbations targeting genes 56--61), IFN-$\gamma$ is consistently the most challenging context, with the highest I-MAE and the lowest DES/PDS. Control and Co-culture are substantially easier; Co-culture is best on I-MAE and PDS, while Control is marginally best on DES. Overall, this points to genuine context-dependent biological differences (rather than a single metric driving the conclusion).

\subsubsection{Correlation between metrics}
To test whether the three metrics are redundant, we compute pairwise correlations across all runs
(and all contexts). Note that I-MAE is a loss (lower is better), while DES and PDS are scores (higher
is better), so negative correlations between I-MAE and the other metrics are expected.

\begin{table}[H]
\centering
\caption{Correlations between metrics over all runs and contexts.}
\label{tab:metric-corr}
\begin{tabular}{lcc}
\toprule
 & Pearson $r$ & Spearman $\rho$ \\
\midrule
I-MAE vs DES & -0.510 & -0.712 \\
I-MAE vs PDS & -0.234 & -0.589 \\
DES vs PDS   &  0.531 &  0.684 \\
\bottomrule
\end{tabular}
\end{table}

\textbf{Interpretation.}
First, runs with lower I-MAE tend to have higher DES and PDS (negative correlations), but the
relationship is not proportional: improvements in I-MAE do not translate linearly into improvements
in DE overlap (DES) or signature ranking (PDS).
Second, the fact that Spearman correlations are stronger than Pearson correlations indicates that
\emph{rank agreement} across metrics is more reliable than \emph{value agreement}: a run that ranks
well on one metric typically ranks well on the others, even if the numerical gaps differ.
Finally, DES and PDS are positively correlated but not interchangeable, motivating the qualitative
checks in Section~\ref{sec:qualitative} where we examine failure modes that can look acceptable under
I-MAE but still produce biologically unfaithful signatures.

\subsection{Qualitative analysis and case studies}
\subsubsection{When I-MAE can be misleading}
\label{sec:qualitative}
I-MAE averages absolute errors over all non-target genes, so it can look favorable even when the
predicted perturbation \emph{signature} is biologically incorrect. We observed three practical failure
modes that motivate reporting DES and PDS alongside I-MAE.

\paragraph{(1) Error averaging can hide signature mistakes.}
Because I-MAE aggregates over many genes, errors on a relatively small set of responsive genes can
be diluted by many near-unchanged genes. In this regime, two runs can have similar I-MAE while
differing substantially in whether they capture the correct direction and ranking of expression changes.
PDS is designed to be sensitive to this by comparing the full ranked delta signature to ground truth.

\paragraph{(2) ``Smooth'' predictions can look good under I-MAE.}
A predictor that regresses toward a context-specific mean profile can achieve a decent average error
even if it underestimates perturbation-specific effects (i.e., it predicts changes that are too small).
This can yield acceptable I-MAE while producing weak or non-specific delta signatures, which is
penalized by PDS and can also reduce DES when true DE effects are not recovered.

\paragraph{(3) Correct mean-level reconstruction does not imply correct DE sets.}
A model may match overall expression levels well but still fail to recover which genes are truly
differentially expressed under a perturbation (wrong DE set, or missing key DE genes). DES directly
tests overlap of DE gene sets and can therefore disagree with I-MAE in such cases.

Overall, these cases explain why the correlations between metrics are imperfect (Table~\ref{tab:metric-corr}):
I-MAE, DES, and PDS are related, but each detects distinct aspects of biological faithfulness.

\section{Implementation and Reproducibility}\label{implementation}
This section documents the practical decisions required to reproduce Bicycle training and to
compute I-MAE/DES/PDS consistently across contexts. In perturbation modeling, these details are
not cosmetic: metric values can shift meaningfully with small changes in normalization, control
selection, and pseudobulk construction.

\subsection{Data handling and context isolation}
We treat each biological context (Control, IFN-\(\gamma\), Co-culture) as a separate dataset and train
a separate model per context. This mirrors the protocol in \cite{rohbeck2024bicycle} and avoids
mixing context effects into the learning objective.

Within each context, we ensure that:
\begin{itemize}
    \item Control (non-targeting) cells are explicitly identified and retained for both DE testing
    (DES) and delta-signature computation (PDS).
    \item The perturbation label per cell is consistent across training, evaluation, and metric code.
    \item The held-out set contains perturbations that are unseen during training (90/10 split).
\end{itemize}

\subsection{Compute environment and stability considerations}
Training Bicycle involves optimizing per-cell variational parameters and can exhibit run-to-run
variance. We therefore report results over repeated runs where feasible and summarize performance
with mean \(\pm\) std. We also observed occasional outliers in I-MAE, consistent with the behavior
noted in \cite{rohbeck2024bicycle}. When using Bicycle in practice, a multi-seed aggregation (e.g., reporting the median prediction/score across several independent training seeds) can provide a more robust estimate than a single run.

% -----------------------------------------------------------------------------
\section{Conclusions and Further Research}\label{conclusions}
\subsection{Conclusions}
This thesis evaluated the Bicycle model on three biological contexts derived from
the Frangieh Perturb-CITE-seq data, using a held-out perturbation split as in the original protocol.
While we did not reproduce the exact numerical I-MAE values reported in \cite{rohbeck2024bicycle},
we recovered the same qualitative behavior: performance is context-dependent, and the relative
difficulty ordering is consistent across metrics.

A central empirical takeaway is that single-score evaluation is insufficient. By adding DES and PDS
to the standard I-MAE evaluation, we found that the three metrics are correlated but not redundant:
runs that improve I-MAE often (but not always) improve DES/PDS, and rank agreement is stronger
than linear agreement (Table~\ref{tab:metric-corr}). This matters in practice because I-MAE can look
favorable even when the model fails to recover the correct differential expression \emph{gene set} or
produces overly weak/non-specific perturbation signatures (Section~\ref{sec:qualitative}).

Finally, comparing contexts shows that generalization difficulty is not uniform: IFN-$\gamma$ is the
most challenging setting under all three metrics on the held-out perturbations, whereas Control and
Co-culture are easier. This reinforces the importance of reporting results per context and using
multiple metrics when judging biological faithfulness of perturbation predictions.

\subsection{Limitations}
Several limitations should be considered.
\begin{itemize}
    \item \textbf{Metric dependence:} DES and PDS are not uniquely defined by the data; they require
    explicit design choices, including the DE test (Wilcoxon rank-sum with tie correction), the
    multiple-testing procedure (Benjamini--Hochberg), the FDR threshold (0.05), and the PDS
    similarity/ranking criterion (Manhattan/\(\ell_1\) distance). Alternative but reasonable choices (e.g.,
    different FDR thresholds or distance metrics) can change absolute scores and, in some cases,
    the relative ranking of methods. We therefore emphasize agreement patterns across metrics
    rather than relying on a single headline number.

    \item \textbf{Data dependence:} All experiments were conducted on the processed
    Perturb-CITE-seq subset used in prior work, restricted to \(G=61\) genes and evaluated per
    biological context. This reduced panel simplifies the prediction task and may overstate
    performance compared to genome-wide settings, where modeling gene--gene dependencies,
    sparsity, and sequencing depth variation is more challenging. Consequently, quantitative
    results should be interpreted as evidence within this benchmark rather than as a direct
    estimate of genome-scale performance.

    \item \textbf{Model class:} Bicycle models latent gene states with a steady-state Ornstein--Uhlenbeck
    (OU) process, which is appropriate for perturbations measured after relaxation but less
    suitable when responses exhibit pronounced transient dynamics, delayed effects, or multiple
    phases. In such regimes, the steady-state assumption can bias inferred regulatory effects and
    limit predictive accuracy for time-dependent responses. Extending the framework to explicit
    time-series observations or non-stationary dynamics is therefore a natural direction for
    future work.
\end{itemize}

\subsection{Further research}
We suggest several directions.
\begin{enumerate}
    \item \textbf{Metric robustness studies:} Systematically vary DE testing methods and similarity
    measures in PDS to assess stability of conclusions.
    \item \textbf{Broader baselines:} Compare \textsc{Bicycle} against strong perturbation predictors under the same train/test perturbation split and the same DES/PDS implementation, including \textsc{GEARS} \cite{roohani2024gears}, \textsc{PerturbNet} \cite{yu2025perturbnet}, and \textsc{CaML} \cite{nilforoshan2023zeroshot}. This tests whether the interpretability of an explicit GRN (\(\beta\)) trades off against predictive performance when judged by signature-level metrics beyond I-MAE.
    \item \textbf{Beyond steady state:} Extend Bicycle-style causal modeling to explicitly time-resolved
    perturbation trajectories where available.
\end{enumerate}

\newpage
\bibliographystyle{alpha}
\bibliography{bibliography}
\addcontentsline{toc}{section}{References}

\end{document}
